% ============================================================================
%  Literature review — AI-Enabled Smart EV Charging and Energy Management
%  Build:  xelatex -> bibtex -> xelatex -> xelatex
% ============================================================================
\documentclass[11pt,a4paper]{article}

\usepackage[UTF8]{ctex}
\usepackage{fontspec}
\setmainfont{Times New Roman}
\setsansfont{Arial}
\setmonofont{Consolas}

\usepackage[a4paper,margin=2.4cm]{geometry}
\usepackage{graphicx}
\usepackage{float}
\usepackage{booktabs}
\usepackage{array}
\usepackage{amsmath}
\usepackage{xcolor}
\usepackage{enumitem}
\usepackage{caption}
\usepackage{fancyhdr}
\usepackage[hidelinks]{hyperref}
\usepackage{parskip}
\usepackage[htt]{hyphenat}
\setlength{\emergencystretch}{2.5em}

\captionsetup{font=small,labelfont=bf}
\setlist[itemize]{leftmargin=1.4em,topsep=2pt,itemsep=2pt}
\setlist[enumerate]{leftmargin=1.6em,topsep=2pt,itemsep=2pt}

\pagestyle{fancy}
\fancyhf{}
\fancyhead[L]{\small FYP — Smart EV Charging Platform}
\fancyhead[R]{\small Literature Review}
\fancyfoot[C]{\small\thepage}
\renewcommand{\headrulewidth}{0.4pt}

\newcommand{\code}[1]{\texttt{\small #1}}

\begin{document}

\begin{center}
  {\LARGE\bfseries Literature Review}\\[6pt]
  {\large AI-Enabled Smart Electric Vehicle Charging\\[2pt]
   and Energy Management Platform}\\[8pt]
  {\normalsize Sub-Project 1 companion document}
\end{center}

% ============================================================================
\section{Scope and selection criteria}

This review covers the five strands the project depends on: machine-learning forecasting of EV
charging demand; the conversion of per-session charging records into hourly load profiles; smart
charging against time-of-use tariffs; open charging datasets; and the measured grid impact of
charging behaviour. It is a working review for a systems-engineering project, so each strand closes
with what it changes in the design rather than with a summary of the field.

Every cited work was accepted only after its DOI resolved through Crossref with matching title,
authors, year and venue. Where a publisher blocks automated access, the abstract was taken from the
Crossref or OpenAlex deposit record rather than from a search result. Works whose abstracts could not
be obtained, and one DOI that does not exist, are listed in Section~\ref{sec:gaps}. Datasets
distributed without a DOI are cited through the peer-reviewed paper that introduces them.

% ============================================================================
\section{Forecasting EV charging demand}\label{sec:forecast}

The strand divides into model-family comparisons and one systematised review.

\paragraph{Model families.} Zhu et al.~\cite{Zhu2019ShortTermEVLoadDL} compare four deep
architectures for station-level hourly load and report a gated recurrent unit network as the best
performer on hourly historical data; their framing is that driver behaviour is stochastic while the
fleet's power draw is large, which is the combination that makes the problem hard.
Vishnu et al.~\cite{Vishnu2023ShortTermEVLoadTSMLDL} compare an auto-regressive model, support vector
regression and an LSTM on Adaptive Charging Network data and report the LSTM as best, with
MAE 4\,kW and RMSE 5.9\,kW --- the most concrete published error figure found in this strand, and
the only one trained on the same data family the project uses as a comparison point.
Zhang et al.~\cite{Zhang2022MCCNNTCN} combine a multi-channel CNN with a temporal convolutional
network and splice in meteorological and calendar covariates, reporting improvements over ANN of
14.09\,\%, LSTM 25.13\,\%, CNN-LSTM 27.32\,\% and TCN 4.48\,\% on urban charging-station load from a
Chinese city. Ge et al.~\cite{Ge2020SpatialTemporalEVLoadProfile} take a different route: an improved
random forest with harmony-search hyperparameter tuning, followed by a bottom-up construction of the
spatial-temporal load of EV clusters.

\paragraph{Method taxonomy.} Alaraj et al.~\cite{Alaraj2025MLChargingDemandReview} survey machine-learning
demand forecasting and argue the now-standard position: unmanaged mass adoption risks network
instability, smart charging can mitigate it, but only given an accurate forecast. That ordering ---
forecast first, schedule second --- is the architecture this project implements.

\paragraph{What this strand settles.} Hourly resolution is the resolution the literature reports at,
and calendar covariates are established as useful rather than novel. It also supplies the reference
error figures against which the project's own accuracy can be discussed, with the caveat that
absolute errors are not comparable across datasets of different scale.

% ============================================================================
\section{From session records to hourly load profiles}\label{sec:profiles}

This strand carries the project's central data-engineering assumption: that a table of sessions ---
each with a connection time, a disconnection time and an energy quantity --- can be converted into an
hourly power profile.

Yi and Scoffield~\cite{Yi2018ResidentialLoadProfileGeneration} build a charging-decision model from
historical residential sessions and estimate the distribution of charging duration \emph{conditional
on parking duration} using kernel density estimation, then generate synthetic charging behaviour that
preserves the parking-to-travel dependency. Uimonen and Lehtonen~\cite{Uimonen2020OfficeBuildingOccupancy}
derive office-building charging profiles from measured \emph{charger occupancy} and use them to ask how
many chargers a building needs --- the closest published analogue to this project's workplace,
midday-peaking site. Munkhammar et al.~\cite{Munkhammar2014BernoulliPEVCharging} condense a Monte Carlo
state-of-charge simulation into hour-by-hour Bernoulli distributions of plug-in status and show that
the aggregate across many vehicles is approximately normal, which is the formal basis for turning
per-vehicle plug-in indicators into a fleet profile. Lahariya et al.~\cite{Lahariya2020SyntheticDataGenerator}
model arrivals as a Poisson-like process with connection time made dependent on arrival time through a
Gaussian mixture, and report generated sessions that are statistically indistinguishable from the real
data they were trained on.

von Hoffen~\cite{vonHoffen2016ChargingTransactionAnalytics} makes the complementary point that
transaction records are not self-sufficient: they acquire their analytical value only after enrichment
with background and environmental context.

\paragraph{What this strand settles.} Charging duration should be modelled as a function of occupancy
rather than as a constant rate, which is the assumption this project implements by spreading a
session's metered energy across its \emph{charging} window rather than across a fixed power. It also
establishes that the profile shape for a workplace site is properly derived from occupancy data, and
that enrichment with calendar (and, where available, weather) context is a documented requirement
rather than an optional refinement.

% ============================================================================
\section{Smart charging, tariffs and peak reduction}\label{sec:scheduling}

Chen et al.~\cite{Chen2012ValleyFilling} formulate joint optimal power flow and EV charging for a
fleet with heterogeneous charging-rate limits and characterise the optimal offline schedule as a
\emph{valley-filling} profile, then give both a cheaper centralised algorithm and a decentralised
online algorithm that tracks that profile. This is the theoretical benchmark against which a
scheduling heuristic can be judged approximately optimal. Ramadan et
al.~\cite{Ramadan2018PeakShavingValleyFilling} solve the related peak-shaving and valley-filling problem
with particle swarm optimisation, allowing vehicle-to-grid discharge during peaks.

On the tariff side, Davis and Bradley~\cite{Davis2012TOURatesEfficacy} test whether time-of-use rates
actually shift charging and find a qualification worth carrying into any scheduling claim: for
plug-in hybrids, the additional fuelling cost incurred by charging off-peak can reduce or cancel the
economic incentive the tariff is meant to provide. Tariff response is therefore not free, and a
cost-saving result that assumes it is free overstates the benefit.

Two studies quantify what managed charging is worth at scale. Szinai et
al.~\cite{Szinai2020ReducedGridOperatingCosts} link mobility and grid-dispatch models for California
and report that 0.95--5 million managed vehicles avoid USD 120--690 million in annual grid
operating costs (up to 10\,\%) and cut renewable curtailment by up to 40\,\% --- but also that
overnight time-of-use charging \emph{increases} curtailment while achieving similar cost savings.
Dixon et al.~\cite{Dixon2020MinimiseEmissionsWindCurtailment} schedule charging against grid carbon
intensity and wind availability, reducing average charging emissions from 35--56 to
28--40\,gCO\textsubscript{2}/km, and estimate that 500{,}000 vehicles could absorb about three
quarters of the curtailment at Scotland's largest onshore wind farm.

\paragraph{What this strand settles.} Valley-filling is the right target shape for a scheduler, and
cost and curtailment are separately measurable objectives that can move in opposite directions. The
project's renewable-alignment objective therefore needs to state which of the two it optimises.

% ============================================================================
\section{Open datasets}\label{sec:datasets}

Amara-Ouali et al.~\cite{AmaraOuali2021EVLoadOpenDataReview} survey open data for reproducible EV
charging-load research and report a search across more than 860 repositories yielding roughly
60 datasets, motivating the survey with the observation that the field commonly lacks
charging-station data of realistic consistency. That finding frames this project's dataset choice:
two municipal session-level releases were selected from that landscape on the criteria of official
provenance, a public-domain licence, metered energy per session, and undisputed download.

Li et al.~\cite{Li2025UrbanEV} publish UrbanEV, an open benchmark covering charging space
availability and electricity consumption in Shenzhen at hourly resolution over six months for more
than 20{,}000 stations, including occupancy, duration, volume and price, with weather and spatial
covariates --- the closest existing precedent for the hourly profile work this project performs.
Lee et al.~\cite{Lee2021ACNSim} present ACN-Sim, whose contribution here is double: it is the
DOI-bearing citation for the Adaptive Charging Network dataset, which has no archived DOI of its own,
and it demonstrates how a session dataset becomes an evaluable charging scenario.
Sadeghianpourhamami et al.~\cite{Sadeghianpourhamami2018EVFlexibility} analyse EV flexibility from
measured charging data; its metadata verified, but no abstract could be obtained, so no substantive
claim is attributed to it here.

% ============================================================================
\section{Charging behaviour and grid impact}\label{sec:impact}

Powell et al.~\cite{Powell2022ChargingInfrastructureGridImpacts} model charging demand across the
United States Western Interconnection and evaluate it against a 2035 economic dispatch. Peak net
electricity demand rises by up to 25\,\% under forecast adoption and by 50\,\% in a
full-electrification stress test. Two of their findings bear directly on this project. First,
locally optimised charging combined with a high share of home charging can \emph{strain} the grid
rather than relieve it. Second, shifting to \emph{uncontrolled daytime} charging reduces storage
requirements, excess non-fossil generation, ramping and emissions. The project's observed profile
already peaks at midday, which is favourable under that finding, but a scheduler that instead
maximises overnight load --- the framing inherited from the project brief --- moves in the opposite
direction. This tension is a finding to be examined, not a detail to be smoothed over.
Verma et al.~\cite{Verma2021ReviewImpactsDistributionNetwork} review distribution-network impacts
across penetration levels, supporting the project's treatment of impact as a function of how many
vehicles charge rather than as a fixed property of EVs.

% ============================================================================
\section{Implications for the project}\label{sec:implications}

\begin{enumerate}
  \item \textbf{Resolution and evaluation design.} One-hour resolution matches the published
        forecasting work, and a rolling-origin out-of-sample protocol with explicit baseline models
        is the minimum needed for the reported accuracy to mean anything. Published error figures
        (\cite{Vishnu2023ShortTermEVLoadTSMLDL}: MAE 4\,kW) come from other datasets and larger sites,
        so they are context, not a target.
  \item \textbf{The session-to-hourly derivation.} Spreading metered energy across a session's
        charging window follows the occupancy-conditional modelling of
        \cite{Yi2018ResidentialLoadProfileGeneration} and \cite{Uimonen2020OfficeBuildingOccupancy}.
        The residual weakness --- no within-session power taper --- is a limitation of the source
        data, and it is stated in the pipeline's output metadata.
  \item \textbf{Enrichment.} \cite{vonHoffen2016ChargingTransactionAnalytics} and
        \cite{Zhang2022MCCNNTCN} both indicate that calendar and weather context raise forecast
        quality. The pipeline already carries calendar covariates; adding weather needs an external
        source and remains open.
  \item \textbf{Baselines.} The reviewed forecasting papers compare model families against each
        other. This project additionally scores three naive baselines in the same table, because a
        learned model that does not beat persistence has not demonstrated value regardless of its
        $R^2$.
  \item \textbf{Scheduling objective.} Valley filling is the characterisation of the optimal
        schedule \cite{Chen2012ValleyFilling}; cost and curtailment are distinct objectives
        \cite{Szinai2020ReducedGridOperatingCosts}; and carbon-intensity-aware scheduling has
        published magnitudes \cite{Dixon2020MinimiseEmissionsWindCurtailment}. The project must
        state which objective it optimises and report the other two as secondary.
  \item \textbf{The daytime-charging tension.} \cite{Powell2022ChargingInfrastructureGridImpacts}
        argues that daytime charging reduces system cost while local optimisation with high home
        charging can strain the grid. Since both of this project's sites already peak at midday,
        the scheduling chapter should test valley filling \emph{and} midday alignment rather than
        assuming the brief's overnight framing.
  \item \textbf{Granularity of the forecast.} The Bernoulli-disaggregation and joint
        arrival-departure models \cite{Munkhammar2014BernoulliPEVCharging,Lahariya2020SyntheticDataGenerator}
        both presume per-vehicle records. One of this project's two sites carries a vehicle
        identifier and the other does not, which is why the aggregate-versus-per-EV question is a
        scope decision rather than a data limitation.
\end{enumerate}

% ============================================================================
\section{Works excluded and gaps}\label{sec:gaps}

Three candidate references were dropped because no abstract could be obtained from any deposit
record, which would have left their contribution unverifiable:
\code{10.1016/j.apenergy.2024.125174}, \code{10.1038/s41560-018-0136-x} and
\code{10.1016/j.rser.2023.113873}. One DOI found during the search does not resolve at all
(\code{10.1109/isgt.2019.8791267}, HTTP 404 at both Crossref and the resolver) and is recorded here so
that it is not reused in good faith. A further DOI resolves to a liquid-fuel demand study rather than
an EV study and was rejected on content.

Four datasets are cited through their introducing papers because they carry no DOI of their own: the
Adaptive Charging Network release \cite{Lee2021ACNSim}, UrbanEV \cite{Li2025UrbanEV}, and the two
municipal releases used in this project, which are documented in the repository's data register.

Outside the scope of this review: power-system optimisation beyond the charging interface, battery
degradation, and vehicle-to-grid economics, all of which matter to the wider platform but not to the
forecasting and disaggregation work reviewed here.

\bibliographystyle{unsrt}
\bibliography{references}

\end{document}
